\ssscp{Many competing synonyms in the region}{02-03-02-01}

Many words have competing synonyms,
near synonyms or variants in the region.
For example, \citet[249]{Blust1993} observes that PMP *ka-labaw ‘rat’
and PCEMP *kanzupay ‘rat’ coexist \ili{east} of the Blust line.
\citet{BlustTrussel2020} list PMP *tiduR/*tuduR,
and PMP *qinəp/PCEMP *qenəp,\footnote{
		\citet{BlustTrussel2020}
		give PCEMP *qenəp as a doublet for PMP *qinəp ‘lie down to sleep’.
		However all their supporting forms are for SHWNG
		and \tsc{Oceanic} languages, suggesting that this form is more appropriate
		for PEMP, rather than PCEMP.
		While they list \ili{Kowiai} as a SHWNG language,
		we reclassify it as \tsc{Seram-Tanimbar-Bomberai} (see \crf{ch:STB}).
		That means there is only one reflex of *qenəp found outside of
		EMP languages in their data.}
in which the local glosses in
the many languages represented consistently hover around ‘sleep, lie down’.
\ili{As} an example of these terms in real languages,
\ili{Buru} has \it{bage} ‘sleep,
lie down’, \it{toro-n} ‘sleep soundly’ (from PMP *tuduR),
while the closely related \ili{Lisela} language on \ili{Buru} has \it{ine}
‘sleep, lie down’ (from PMP *qinəp).
There are additional forms listed in Blust
and Trussel for \tsc{PAn}
and PMP under ‘sleep’ as well.
It is common for there to be many competing ways to talk about the same thing.
The distribution of these competing terms (and many others) are
in languages in close proximity to each other.
(See the distribution maps throughout this work.)

\ili{As} another of many examples,
we have already noted that ‘banana’ has multiple generic terms
scattered through the region.
Reflexes of PMP *punti coexist alongside {\#}muku from Papuan
sources (\srf{02-03-01-07}).
\citet{DonohueDenham2009} document additional reflexes also extant
in the region, that all mean ‘banana’.
These are also found outside of our target region,
and must all be considered as competing synonyms for ‘banana’.
Data in \trf{02-06} are gleaned from \citet{DonohueDenham2009}; with our subgroups
and some corrections added.
The reconstructed forms are also modified in some cases
to better account for the reflexes observed in the region.

\begin{table}
	\centering\tcp{Reflexes of competing terms for banana}{02-06}
\begin{threeparttable}
	\begin{tabular}{lll}\lsptoprule
Language	&	\hp{*}‘banana’	&	Subgroup		\\ \midrule
	&	\hp{*}*kalauR	&				\\ \midrule
\ili{Bima}	&	\hp{**}kalo	&	\tsc{\ili{Bima}-Lembata}			\\
\ili{Komodo}	&	\hp{**}kalo	&	\tsc{\ili{Bima}-Lembata}				\\
\ili{Kambera}	&	\hp{**}kaluu	&	\tsc{\ili{Bima}-Lembata}				\\
\ili{Wanukaka}	&	\hp{**}kaloʔa	&	\tsc{\ili{Bima}-Lembata}				\\
\ili{Barakai}	&	\hp{**}kalaur	&	\tsc{Aru}			\\
SW \ili{Tarangan}	&	\hp{**}kalor	&	\tsc{Aru}				\\
	&	**kula	&					\\
\ili{Asilulu}	&	\hp{**}kula	&	\tsc{\ili{Ambon}-Seram}			\\
\ili{Laha}	&	\hp{**}kula	&	\tsc{\ili{Ambon}-Seram}				\\
\ili{Haruku}	&	\hp{**}kura	&	\tsc{\ili{Ambon}-Seram}				\\
\ili{Saparua}	&	\hp{**}kula	&	\tsc{\ili{Ambon}-Seram}				\\
\ili{Sou Amana Teru}	&	\hp{**}kula	&	\tsc{\ili{Ambon}-Seram}				\\ \midrule
	&	*telewa	&			\mc{1}{r}{also in SHWNG}		\\ \midrule
\ili{Nusalaut}	&	\hp{*}telewa	&	\tsc{\ili{Ambon}-Seram}				\\
\ili{Teluti} (\ili{Tehoru})	&	\hp{*}telewa	&	\tsc{\ili{Ambon}-Seram}				\\
East \ili{Tarangan}	&	\hp{*}taragwar	&	\tsc{Aru}				\\ \midrule
	&	*biak	&			\mc{1}{r}{also in SHWNG and Papuan lgs.}	\\ \midrule
\ili{Teun}	&	\hp{*}fiwa	&	\tsc{Timor-Babar}				\\
\ili{Nila}	&	\hp{*}hia	&	\tsc{Timor-Babar}				\\
\ili{Ambelau}	&	\hp{*}bia-e	&	\tsc{Sula-Buru}				\\
\ili{Lisela}	&	\hp{*}fia-t\su{†}	&	\tsc{Sula-Buru}				\\
\ili{Sanana}	&	\hp{*}fia	&	\tsc{Sula-Buru}				\\
\ili{Mangole}	&	\hp{*}fia	&	\tsc{Sula-Buru}				\\
\ili{Kadai}\su{‡} &	\hp{*}fia	&	\tsc{Taliabu} (\tsc{Celebic}) 			\\
\ili{Taliabu}	&	\hp{*}fia	&	\tsc{Taliabu} (\tsc{Celebic}) 			\\
\ili{Soboyo}	&	\hp{*}fiak	&	\tsc{Taliabu} (\tsc{Celebic}) 			\\
		\lspbottomrule
	\end{tabular}
			\begin{tablenotes}\tnsize
				\item[†]	{Replacing \citet{DonohueDenham2009} listing \ili{Buru} \it{fua-t}.}
				\item[‡]	{Corrected from Tai-\ili{Kadai}.}
			\end{tablenotes}
	\end{threeparttable}
\end{table}